\documentclass[10pt,a4paper]{amsart}
\usepackage[margin=19mm]{geometry}
\usepackage{amsmath,amssymb,mathtools}
\usepackage[scr=rsfs,cal=euler]{mathalfa}
\usepackage{xfrac}
\usepackage[unicode,hidelinks,bookmarks=false]{hyperref}
\newcommand{\ie}{i.e.\ }
\newcommand{\cf}{cf.\ }
\newcommand{\resp}{respectively\ }
\newcommand{\Subsection}{\subsection}
\providecommand{\nobreakdash}{\nobreak\hbox{-}}
\setlength{\emergencystretch}{2em}
\multlinegap0pt
\usepackage{bm}
\usepackage{bbm}
\usepackage{dsfont}
\usepackage{xspace}
\usepackage{cancel}
\usepackage{mathtools}
\usepackage{amsthm}
\makeatletter
\@ifundefined{theorem}{\newtheorem{theorem}{Theorem}}{}
\makeatother
\title[The $C^*$-algebras of locally finite undirected graphs: A co]{The $C^*$-algebras of locally finite undirected graphs: A complete description of their K-theory}
\author{D. Pask}
\date{Source: arXiv:2607.08271v1 (CC BY 4.0)}
\begin{document}
\maketitle

\noindent\emph{Source and licence.} The $C^*$-algebras of locally finite undirected graphs: A complete description of their K-theory. Version arXiv:2607.08271v1 (CC BY 4.0), distributed under the Creative Commons Attribution 4.0 International licence, as recorded in the arXiv licence metadata for this exact version and shown on the abstract page https://arxiv.org/abs/2607.08271v1. Full rights notice: licenses/arxiv-2607.08271-CC-BY-4.0.txt.\par\smallskip

\noindent\textit{Editorial scope.} Counted unit: the Theorem whose source label is \texttt{thm:general-realisation} in the article (source0.tex lines 284--286, source label \texttt{thm:general-realisation}), delivered together with the article's own complete original proof of it (source0.tex lines 288--312, one \texttt{proof} environment of 25 lines). No mathematics is written, repaired or re-derived: statement and proof are the article's own text line for line. Required context retained uncounted: none. Every definition, display and label of those lines is retained unchanged. Omitted from this package and not redistributed: 1--283, 287--287, 313--1088. Nothing inside a retained excerpt is omitted or altered: every retained body is delivered exactly as the article's lines are written, so no omission receipt of any kind accompanies this package. Citations: the retained excerpts carry no citation command at all, so no bibliography entry of the article is transcribed and no reference is redistributed. Reference adaptation: none was needed and none was made; every reference inside the delivered unit resolves inside it, so no unresolved reference or citation is produced. Printed slips kept literal: no printed word, symbol, hypothesis or number of the counted statement or of its proof was altered. Layout and class: the article's own class is \texttt{amsart} with packages including amsmath, amssymb, amsthm, mathrsfs, multirow, tikz, geometry, hyperref; the wrapper is the lane's pinned \texttt{amsart} editorial wrapper at 10pt on A4, which restates the theorem-like environments the retained text needs and adds the article's own macro definitions that the retained text needs (none), with extra wrapper packages (bm, bbm, dsfont, xspace, cancel, mathtools). The delivered numbering is the unit's own. Assets: no retained span contains a figure environment, a graphics-inclusion command, a PDF-page inclusion, a table or a TikZ picture; no image file is copied into this package, the package declares no asset dependency and no image file is redistributed with it. Licence: version arXiv:2607.08271v1 of this article is distributed under the Creative Commons Attribution 4.0 International licence, as recorded in the arXiv licence metadata for this exact version and shown on the abstract page https://arxiv.org/abs/2607.08271v1. A full rights notice is delivered at licenses/arxiv-2607.08271-CC-BY-4.0.txt. Characters: every retained excerpt is pure ASCII, so the delivered engine, XeLaTeX with its default Latin Modern text font, reports no missing character.
\begin{theorem}\label{thm:general-realisation}
For every locally finite graph $\Gamma$ we have $C^*(\Gamma)\cong C^*(E_\Gamma)$.
\end{theorem}

\begin{proof}[Sketch of proof]
Let $\{p_u,s_f\}$ be a universal generating family for $C^*(\Gamma)$. For $f\in\Gamma^1\subset
E_\Gamma^0$ put $Q_f=s_fs_f^*$; for each tail vertex put $Q_{u_f}=s_f^*s_f$, the projection left
unresolved by (R2); for $fe\in\Gamma^2$ put $S_{fe}=s_fs_es_e^*$; and for each tail edge put
$T_{w_f}=s_f$. Relations (R1) and (R2) show that $\{Q_v : v\in E_\Gamma^0\}$ is a family of
mutually orthogonal projections. Since $fe\in\Gamma^2$ guarantees that (R2) applies at $f$, we have
$s_es_e^*\le s_f^*s_f$ whenever $fe\in\Gamma^2$, whence
$S_{fe}^*S_{fe}=s_es_e^*=Q_{s_P(fe)}$; and for $f$ with $|r^{-1}(s(f))|\ge 2$, summing over the
edges of $E_\Gamma$ with range $f$,
\[
\sum_{r(e)=s(f),\,e\neq\bar f}S_{fe}S_{fe}^{*}
= s_f\Big(\sum_{r(e)=s(f),\,e\neq\bar f}s_es_e^{*}\Big)s_f^{*}
= s_f\,s_f^{*}s_f\,s_f^{*}=Q_f,
\]
again by (R2). For the tail edges,
$T_{w_f}^*T_{w_f}=s_f^*s_f=Q_{u_f}=Q_{s_P(w_f)}$, and since $w_f$ is the unique edge of $E_\Gamma$
with range $f$, the relation at $f$ is $T_{w_f}T_{w_f}^*=s_fs_f^*=Q_f$; the vertices $u_f$ are
sources and carry no relation. Hence $\{Q_v,S_{fe},T_{w_f}\}$ is a Cuntz--Krieger
$E_\Gamma$-family in $C^*(\Gamma)$, and the universal property of $C^*(E_\Gamma)$ yields a
homomorphism $\pi\colon C^*(E_\Gamma)\to C^*(\Gamma)$ intertwining the gauge actions. Since each
$Q_v\neq 0$, the gauge-invariant uniqueness theorem shows $\pi$ is injective. For surjectivity,
$s_f=T_{w_f}$ when $|r^{-1}(s(f))|=1$, while otherwise (R2) gives
$s_f=s_fs_f^*s_f=\sum_{r(e)=s(f),\,e\neq\bar f}S_{fe}$; the projections $p_u$ then lie in the image
by (R1).
\end{proof}

\end{document}
